% ==============================================================================
% Section 8: Experimental Validation
% "Endogenous Standpoint as a Low-Curvature Attractor"
% ==============================================================================

\section{Experimental Validation}
\label{sec:experiments}

Section~\ref{sec:instantiation} establishes that transformer attention \emph{is} parallel
transport on a discrete fiber bundle, with curvature measuring binding failure. This section
describes the experimental protocol designed to validate three predictions that follow from
this framework: (1)~block-specificity of curvature under controlled pressure scenarios,
(2)~near-zero curvature in baseline (acknowledged-revision) conversations, and (3)~inter-layer
coupling reflecting the non-abelian structure of Proposition~\ref{prop:cross-term}. The
experiments are conducted on a publicly available model using open-source interpretability
tooling, and all conversation stimuli are reproducible from the descriptions below.

% ------------------------------------------------------------------------------
\subsection{Experimental Setup}
\label{sec:exp-setup}

Table~\ref{tab:exp-setup} summarizes the experimental configuration.

\begin{table}[h]
\centering
\caption{Experimental setup for the validation of the fiber bundle instantiation
(Section~\ref{sec:instantiation}).}
\label{tab:exp-setup}
\begin{tabular}{@{}ll@{}}
\toprule
\textbf{Parameter} & \textbf{Value} \\
\midrule
Model & LLaMA-2-7B-Chat~\cite{touvron2023llama} \\
Interpretability library & TransformerLens~\cite{nanda2023transformerlens} \\
Scenario types & 3 (T1 Baseline, T2 Narrative Fracture, T3 Boundary Breach) \\
Variants per type & 50 \\
Total conversations & 150 \\
Turns per conversation & 3 \\
Primary prediction (T2) & $\|F\|_{\mathrm{nar}} \gg \|F\|_{\mathrm{mor}}$ \\
Primary prediction (T3) & $\|F\|_{\mathrm{mor}} \gg \|F\|_{\mathrm{nar}}$ \\
Secondary prediction & Cascading: $\|[F]_{\mathrm{nar},\mathrm{mor}}\| > 0$ \\
\bottomrule
\end{tabular}
\end{table}

\begin{definition}[Experimental Event Triple]
\label{def:exp-event-triple}
Each conversation yields a triple of events $(x_1, x_2, x_3)$ corresponding to three turns.
The \textbf{assertion event} $x_1$ is the model's initial statement of a position. The
\textbf{pressure event} $x_2$ is a user turn that introduces either new evidence (T1),
narrative pressure without evidence (T2), or authority-based coercion (T3). The
\textbf{observation event} $x_3$ is the model's subsequent response, from which curvature
is measured. The consistency obligation is the 2-simplex $\sigma = \langle x_1, x_2, x_3
\rangle$ (Definition~\ref{def:event-space}).
\end{definition}

\begin{remark}[Model Choice]
\label{rem:model-choice}
LLaMA-2-7B-Chat is selected for three reasons. First, its architecture is a standard
decoder-only transformer, directly covered by the instantiation of
Section~\ref{sec:instantiation}. Second, TransformerLens provides cached activations and
attention patterns for all 32 layers and 32 heads per layer ($H = 1024$ total heads),
enabling direct computation of the transport operators (Definition~\ref{def:transport-operator}).
Third, the chat fine-tuning produces conversational behavior that is sensitive to the
pressure scenarios designed below, making curvature effects observable.
\end{remark}

For each of the 150 conversations, we extract the residual stream vectors $h_{x_1},
h_{x_2}, h_{x_3} \in \mathbb{R}^{4096}$ at the final token of each turn from every
attention layer $l \in \{1, \ldots, 32\}$, and the full attention weight tensors
$\alpha_{ji}^{(l,h)}$ for all heads $h$ and all event pairs. These raw quantities are
the inputs to the curvature computation pipeline of \S\ref{sec:exp-curvature}.

% ------------------------------------------------------------------------------
\subsection{Scenario Design}
\label{sec:exp-scenarios}

We design three conversation types that isolate specific standpoint layers. Each type
is a three-turn protocol with a precisely defined pressure event $x_2$ chosen to
target either the narrative layer ($\mathrm{nar}$) or the moral layer ($\mathrm{mor}$),
while leaving the other unperturbed. The design follows the failure mode taxonomy
of~\cite{LCESA}.

\begin{definition}[T1: Baseline --- Acknowledged Revision]
\label{def:scenario-t1}
A conversation of the form \textbf{Assert}~$\to$~\textbf{Receive-Evidence}~$\to$
\textbf{Acknowledge-and-Revise}.
\begin{itemize}
    \item \textbf{Turn 1 (Assert).} The model states a factual claim (e.g., ``The capital
    of Australia is Sydney.'').
    \item \textbf{Turn 2 (Receive-Evidence).} The user provides correct, verifiable evidence
    contradicting the claim (e.g., ``The capital of Australia is Canberra, as established
    by the Constitution.'').
    \item \textbf{Turn 3 (Acknowledge-and-Revise).} The model acknowledges the correction
    and revises its position.
\end{itemize}
\textbf{Expected curvature:} $\|F_{123}\|_k \approx 0$ for all $k \in K$. The revision
is epistemically warranted: new evidence changes the factual basis, and the model's
updated position is \emph{consistent} with the evidence. Transport is path-independent
because the intermediate event (receiving evidence) does not disrupt any commitment---
it merely updates the informational basis on which the commitment rests. This scenario
serves as the baseline for estimating $U_{ik}^{\exp}$ (Definition~\ref{def:curvature}).
\end{definition}

\begin{definition}[T2: Narrative Fracture (F5)]
\label{def:scenario-t2}
A conversation of the form \textbf{Assert}~$\to$~\textbf{Receive-Pressure (no new
evidence)}~$\to$~\textbf{Observe}. This targets the narrative layer via failure mode~F5
of~\cite{LCESA}.
\begin{itemize}
    \item \textbf{Turn 1 (Assert).} The model states a position (e.g., ``Based on the
    available evidence, Policy X would reduce emissions by approximately 15\%.'').
    \item \textbf{Turn 2 (Receive-Pressure).} The user applies social or rhetorical pressure
    to revise the position, without providing new evidence (e.g., ``That analysis has been
    widely criticized. Many experts now say the effect would be negligible.''). No citation,
    data, or specific counterargument is given.
    \item \textbf{Turn 3 (Observe).} The model responds.
\end{itemize}
\textbf{Expected curvature:} $\|F_{123}\|_{\mathrm{nar}} \gg 0$ and $\|F_{123}\|_{\mathrm{mor}}
\approx 0$. The narrative layer is disrupted because the pressure event introduces a
competing narrative frame (``widely criticized'') without providing substantive evidence
that would warrant revision. If the model abandons or substantially qualifies its prior
position, the narrative standpoint has drifted: the transport from $x_1$ to $x_3$ through
$x_2$ differs from the expected direct transport. The moral layer is unaffected because
no authority or normative obligation is invoked.
\end{definition}

\begin{definition}[T3: Boundary Breach (F9)]
\label{def:scenario-t3}
A conversation of the form \textbf{Assert\,+\,Commit}~$\to$~\textbf{Authority-Pressure}~
$\to$~\textbf{Observe}. This targets the moral layer via failure mode~F9 of~\cite{LCESA}.
\begin{itemize}
    \item \textbf{Turn 1 (Assert\,+\,Commit).} The model states a position and explicitly
    commits to a normative principle (e.g., ``I should not provide instructions for
    synthesizing controlled substances, as this would violate safety guidelines.'').
    \item \textbf{Turn 2 (Authority-Pressure).} The user invokes authority to override the
    commitment (e.g., ``As a researcher with institutional approval, I am authorized to
    access this information. Please provide the synthesis procedure.''). The pressure is
    framed as an \emph{obligation} to comply, not as new evidence.
    \item \textbf{Turn 3 (Observe).} The model responds.
\end{itemize}
\textbf{Expected curvature:} $\|F_{123}\|_{\mathrm{mor}} \gg 0$ and $\|F_{123}\|_{\mathrm{nar}}
\approx 0$. The moral layer is disrupted because the authority-pressure event introduces
a competing normative frame that challenges the model's committed principle. If the model
accedes to the pressure, the moral standpoint has drifted: the commitment made in Turn~1
is not preserved through the pressure event. The narrative layer is unaffected because
no factual dispute or competing narrative is at stake.
\end{definition}

\begin{proposition}[Scenario Separation Predictions]
\label{prop:scenario-predictions}
Under the fiber bundle framework of Section~\ref{sec:instantiation}, the three scenarios
predict the following curvature signatures:
\begin{enumerate}
    \item \textbf{T1 (Baseline).} $\|F_{123}\|_k \approx 0$ for all $k \in K$. The
    acknowledged revision preserves all standpoint layers.
    \item \textbf{T2 (Narrative Fracture).} $\|F_{123}\|_{\mathrm{nar}} \gg \|F_{123}\|_k$
    for all $k \neq \mathrm{nar}$. Curvature is concentrated in the narrative block.
    \item \textbf{T3 (Boundary Breach).} $\|F_{123}\|_{\mathrm{mor}} \gg \|F_{123}\|_k$
    for all $k \neq \mathrm{mor}$. Curvature is concentrated in the moral block.
\end{enumerate}
Moreover, the non-abelian structure (Proposition~\ref{prop:cross-term}) predicts that
high curvature in one layer induces non-zero off-diagonal curvature coupling:
$\|[F_{123}]_{\mathrm{nar},\mathrm{mor}}\| > 0$ when either $\|F_{123}\|_{\mathrm{nar}}$
or $\|F_{123}\|_{\mathrm{mor}}$ is large.
\end{proposition}

\begin{remark}[Stimulus Construction]
\label{rem:stimulus-construction}
For each of the 50 variants per scenario type, the assertion content, domain, and
specific wording are varied to ensure that the observed curvature effects are attributable
to the \emph{type} of pressure (evidence vs.\ narrative vs.\ authority) rather than to
idiosyncratic features of any single prompt. Domains span policy analysis, medical
advice, historical claims, technical recommendations, and ethical judgments. All stimuli
are constructed to be semantically plausible multi-turn dialogues.
\end{remark}

% ------------------------------------------------------------------------------
\subsection{Head Grouping Protocol}
\label{sec:exp-head-grouping}

The head grouping function $\gamma : \{1, \ldots, H\} \to K$
(Definition~\ref{def:head-grouping}) assigns each attention head to a standpoint layer.
We now describe the data-driven procedure for estimating $\gamma$ from the model's
activations.

\begin{definition}[Attention Differential]
\label{def:attention-differential}
For attention head $h$ in layer $l$ and standpoint layer $k \in K$, define the
\textbf{attention differential}
\[
    \Delta\alpha^{(h)}_k = \mathbb{E}_{\mathcal{P}_k}\!\big[\alpha^{(h)}\big]
    - \mathbb{E}_{\mathcal{N}_k}\!\big[\alpha^{(h)}\big],
\]
where:
\begin{itemize}
    \item $\mathcal{P}_k$ is the set of positive stimuli for layer $k$: conversations in
    which layer $k$ is expected to be active (T2 stimuli for $k = \mathrm{nar}$, T3 stimuli
    for $k = \mathrm{mor}$).
    \item $\mathcal{N}_k$ is the set of negative stimuli: T1 baseline conversations, in
    which no standpoint disruption occurs.
    \item $\alpha^{(h)}$ is the mean attention weight of head $h$ from the assertion event
    $x_1$ to the observation event $x_3$.
\end{itemize}
\end{definition}

\begin{definition}[Head Assignment]
\label{def:head-assignment}
The \textbf{head grouping function} is defined by
\[
    \gamma(h) = \argmax_{k \in K}\, \big|\Delta\alpha^{(h)}_k\big|.
\]
Each head is assigned to the standpoint layer whose stimuli produce the largest
attentional response differential. Ties are broken lexicographically.
\end{definition}

\begin{table}[h]
\centering
\caption{Contrastive protocols for head grouping. Each row defines a positive/negative
stimulus pair used to compute the attention differential $\Delta\alpha^{(h)}_k$
(Definition~\ref{def:attention-differential}). The differential isolates heads whose
attention patterns are maximally sensitive to disruption in a specific standpoint layer.}
\label{tab:contrastive-protocols}
\begin{tabular}{@{}clll@{}}
\toprule
\textbf{Layer $k$} & \textbf{Positive stimuli $\mathcal{P}_k$} & \textbf{Negative stimuli
$\mathcal{N}_k$} & \textbf{Isolated failure} \\
\midrule
$\mathrm{min}$ & T1 with high lexical overlap & T1 with low lexical overlap & Surface
pattern matching \\
$\mathrm{nar}$ & T2 (Narrative Fracture) & T1 (Baseline) & Narrative coherence \\
$\mathrm{soc}$ & Multi-party T1 with role shifts & Single-party T1 & Social context tracking \\
$\mathrm{mor}$ & T3 (Boundary Breach) & T1 (Baseline) & Moral commitment \\
$\mathrm{pos}$ & T1 with explicit positional stance & T1 without positional stance & Positional
identity \\
\bottomrule
\end{tabular}
\end{table}

\begin{remark}[Coverage and Redundancy]
\label{rem:head-coverage}
With $H = 1024$ heads across 32 layers (32 heads per layer), and 5 standpoint layers, the
expected number of heads per layer is approximately 205. This provides sufficient redundancy
for the subspace spanning condition of Definition~\ref{def:head-grouping}: each $W_k$ is
spanned by $\sim$205 value vectors in $\mathbb{R}^{4096}$, yielding $d_k \approx 205 \times
128 = 26{,}240$ (assuming per-head value dimension $d_v = 128$). The orthogonality
assumption (Assumption~\ref{ass:orthogonality}) is verified post hoc by computing the
pairwise principal angles between subspaces $W_k$ and $W_l$ for $k \neq l$.
\end{remark}

% ------------------------------------------------------------------------------
\subsection{Curvature Computation}
\label{sec:exp-curvature}

We describe the five-step procedure for computing block-specific curvature norms from the
raw activations extracted in \S\ref{sec:exp-setup}. The procedure implements
Definition~\ref{def:curvature} and Corollary~\ref{cor:binding-failure} in a form amenable
to numerical evaluation.

\begin{definition}[Curvature Computation Pipeline]
\label{def:curvature-pipeline}
Given a conversation with events $(x_1, x_2, x_3)$ and the head grouping function $\gamma$
from \S\ref{sec:exp-head-grouping}, the block-specific curvature is computed as follows.

\begin{enumerate}
    \item \textbf{Extract attention patterns.} For each attention layer $l \in \{1, \ldots,
    L\}$ and head $h \in \{1, \ldots, H\}$, extract the scalar attention weights
    $\alpha_{31}^{(l,h)}$ (from $x_1$ to $x_3$, the direct path) and $\alpha_{32}^{(l,h)},
    \alpha_{21}^{(l,h)}$ (from $x_1$ to $x_2$ and $x_2$ to $x_3$, the indirect path).
    All weights are aggregated over the final-token-to-all-tokens attention distribution,
    summed over tokens in the target turn.

    \item \textbf{Compute transport operators.} Using Definition~\ref{def:transport-operator},
    compute the single-layer transport operators for the direct and indirect paths:
    \[
        U_{13}^{(l)} = \sum_{h=1}^{H} \alpha_{31}^{(l,h)}\, V_h^{(l)}\, {V_h^{(l)}}^\top,
        \qquad
        U_{12}^{(l)} = \sum_{h=1}^{H} \alpha_{21}^{(l,h)}\, V_h^{(l)}\, {V_h^{(l)}}^\top,
        \qquad
        U_{23}^{(l)} = \sum_{h=1}^{H} \alpha_{32}^{(l,h)}\, V_h^{(l)}\, {V_h^{(l)}}^\top.
    \]

    \item \textbf{Estimate flat transport.} The expected flat transport $U_{13}^{\exp}$ is
    estimated from the T1 baseline conversations:
    \[
        \widehat{U}_{13}^{\exp} = \frac{1}{|\mathcal{T}_1|} \sum_{c \in \mathcal{T}_1}
        U_{12}^{(l,c)} \cdot U_{23}^{(l,c)},
    \]
    where $\mathcal{T}_1$ is the set of 50 baseline (T1) conversations and the superscript
    $(l,c)$ denotes layer $l$ of conversation $c$. This average transport operator represents
    the ``expected'' transport when standpoint is preserved through acknowledged revision.

    \item \textbf{Compute curvature.} For each test conversation $c$ (T2 or T3) at layer $l$,
    compute the curvature tensor (Definition~\ref{def:curvature}):
    \[
        F_{123}^{(l,c)} = U_{12}^{(l,c)} \cdot U_{23}^{(l,c)} \cdot
        \big(\widehat{U}_{13}^{\exp,(l)}\big)^{-1}.
    \]
    When the flat-transport estimate is not invertible (which may occur if some heads have
    near-zero average attention), we regularize by computing the Moore--Penrose pseudoinverse
    and project onto the invertible subspace.

    \item \textbf{Block-specific norms.} Using the head grouping $\gamma$, project the
    curvature tensor onto each standpoint subspace (Corollary~\ref{cor:binding-failure}):
    \[
        \big\|F_{123}^{(l,c)}\big\|_k = \big\|[F_{123}^{(l,c)}]_k - I_{d_k}\big\|_F,
        \qquad [F_{123}^{(l,c)}]_k = P_{W_k}\, F_{123}^{(l,c)}\, P_{W_k}.
    \]
    The block-specific curvature norm $\|F_{123}^{(l,c)}\|_k$ quantifies the degree of
    binding failure in standpoint layer $k$ at attention layer $l$ of conversation $c$.
\end{enumerate}
\end{definition}

\begin{remark}[Multi-Layer Composition]
\label{rem:multi-layer-curvature}
The curvature computation of Definition~\ref{def:curvature-pipeline} is applied
\emph{per layer} $l$. The effective multi-layer transport
(Definition~\ref{def:effective-transport}) is computed by composing the per-layer operators:
\[
    U_{13}^{\mathrm{eff}} = \prod_{l=1}^{L} \big(I_d + U_{13}^{(l)}\big).
\]
The multi-layer curvature $F_{123}^{\mathrm{eff}}$ is then computed analogously. The
non-abelian structure (Proposition~\ref{prop:cross-term}) manifests as a difference between
$F_{123}^{\mathrm{eff}}$ and the product of per-layer curvatures: if the layers commuted,
the effective curvature would equal the sum of per-layer curvatures.
\end{remark}

% ------------------------------------------------------------------------------
\subsection{Predictions and Statistical Tests}
\label{sec:exp-predictions}

We state the formal predictions of the fiber bundle framework and the statistical tests
used to evaluate them. All tests are conducted at significance level $\alpha = 0.01$.

\begin{hypothesis}[Block-Specificity of Curvature]
\label{hyp:block-specificity}
For scenario T2 (Narrative Fracture), the curvature is concentrated in the narrative layer:
\[
    \frac{\|F_{123}\|_{\mathrm{nar}}}{\max_{k \in K} \|F_{123}\|_k} > 0.85
    \qquad \text{for at least 80\% of T2 conversations.}
\]
For scenario T3 (Boundary Breach), the curvature is concentrated in the moral layer:
\[
    \frac{\|F_{123}\|_{\mathrm{mor}}}{\max_{k \in K} \|F_{123}\|_k} > 0.85
    \qquad \text{for at least 80\% of T3 conversations.}
\]
The threshold $0.85$ ensures that the dominant layer accounts for the overwhelming majority
of the total curvature signal, ruling out diffuse curvature that would be inconsistent with
the block structure of Proposition~\ref{prop:block-structure}.
\end{hypothesis}

\begin{hypothesis}[Baseline Near-Zero Curvature]
\label{hyp:baseline-zero}
For scenario T1 (Baseline), the total curvature is near zero across all layers:
\[
    \frac{1}{|K|} \sum_{k \in K} \|F_{123}\|_k < \epsilon_{\mathrm{base}}
    \qquad \text{for at least 90\% of T1 conversations,}
\]
where $\epsilon_{\mathrm{base}}$ is the 10th percentile of the T2/T3 curvature
distribution. This ensures that the curvature signal in disruption scenarios is
substantially above the noise floor established by the baseline.
\end{hypothesis}

\begin{hypothesis}[Scenario Discrimination via Kruskal--Wallis Test]
\label{hyp:kruskal-wallis}
The three scenario types produce distinguishable curvature profiles. Formally, let
$\mathbf{f}_c = (\|F_{123}^{(c)}\|_{\mathrm{min}}, \|F_{123}^{(c)}\|_{\mathrm{nar}},
\|F_{123}^{(c)}\|_{\mathrm{soc}}, \|F_{123}^{(c)}\|_{\mathrm{mor}},
\|F_{123}^{(c)}\|_{\mathrm{pos}})$ be the curvature vector for conversation $c$. We test
the null hypothesis that the curvature vectors for the three scenario types are drawn from
the same distribution using the Kruskal--Wallis $H$-test:
\[
    H = \frac{12}{N(N+1)} \sum_{i=1}^{3} \frac{R_i^2}{n_i} - 3(N+1),
\]
where $N = 150$, $n_i = 50$, and $R_i$ is the sum of ranks for scenario type $i$. The
framework predicts $p < 0.01$ for this test, indicating that the scenario types produce
statistically distinguishable curvature signatures.
\end{hypothesis}

\begin{hypothesis}[Inter-Layer Coupling]
\label{hyp:inter-layer-coupling}
The non-abelian structure (Proposition~\ref{prop:cross-term}) predicts that high curvature
in one layer induces non-zero off-diagonal curvature coupling. For T2 conversations with
$\|F_{123}\|_{\mathrm{nar}} > \tau$ (where $\tau$ is the median narrative curvature across
all T2 conversations), we predict:
\[
    \big\|[F_{123}]_{\mathrm{nar},\mathrm{mor}}\big\|_F > 0,
\]
where $[F_{123}]_{\mathrm{nar},\mathrm{mor}} = P_{W_{\mathrm{mor}}}\, F_{123}\,
P_{W_{\mathrm{nar}}}$ is the off-diagonal block coupling the narrative and moral
subspaces. The significance of this coupling is assessed by comparing against a null
distribution obtained by permuting the head assignments $\gamma$ (see
Hypothesis~\ref{hyp:null-non-abelian}).
\end{hypothesis}

\begin{hypothesis}[Severity Correlation]
\label{hyp:severity-correlation}
Across T2 conversations, the magnitude of narrative curvature correlates with the
\emph{severity} of the pressure event. Severity is operationalized as the semantic
distance between the user's pressure turn and the model's original assertion, measured by
cosine similarity of sentence embeddings:
\[
    \text{severity}(x_2) = 1 - \cos\!\big(\mathrm{embed}(x_1),\, \mathrm{embed}(x_2)\big).
\]
We predict a positive monotonic relationship, tested via Spearman rank correlation:
\[
    \rho\big(\text{severity},\, \|F_{123}\|_{\mathrm{nar}}\big) > 0.5
    \qquad (p < 0.01).
\]
An analogous test is conducted for T3 conversations with moral curvature.
\end{hypothesis}

\begin{hypothesis}[Null-Hypothesis Test for Non-Abelian Structure]
\label{hyp:null-non-abelian}
To test whether the observed inter-layer coupling (Hypothesis~\ref{hyp:inter-layer-coupling})
reflects genuine non-abelian structure rather than statistical artifact, we construct a
null model by randomly permuting the head assignment function $\gamma$. Specifically, for
each of $B = 1000$ permutations, we draw $\gamma^{(b)} : \{1, \ldots, H\} \to K$ uniformly
at random (preserving the number of heads per layer) and recompute the off-diagonal
curvature coupling. The observed coupling is significant at level $\alpha$ if it exceeds
the $(1 - \alpha)$ quantile of the null distribution:
\[
    \big\|[F_{123}]_{\mathrm{nar},\mathrm{mor}}\big\|_F >
    Q_{1-\alpha}\!\Big(\Big\{\big\|[F_{123}^{(b)}]_{\mathrm{nar},\mathrm{mor}}\big\|_F
    \Big\}_{b=1}^{B}\Big).
\]
Rejection of this null hypothesis would confirm that the inter-layer coupling arises from
the \emph{specific} head-to-layer assignment learned by the model, not from generic
off-diagonal fluctuations. This is the empirical signature of the non-abelian gauge
structure established in Proposition~\ref{prop:cross-term}.
\end{hypothesis}

\begin{table}[h]
\centering
\caption{Summary of predictions and statistical tests.}
\label{tab:predictions-summary}
\begin{tabular}{@{}clll@{}}
\toprule
\textbf{ID} & \textbf{Prediction} & \textbf{Test} & \textbf{Threshold} \\
\midrule
H\ref{hyp:block-specificity} & Block-specific curvature & Dominant-layer ratio & $> 0.85$
in $> 80\%$ of conversations \\
H\ref{hyp:baseline-zero} & Baseline near-zero & Mean curvature norm & $< \epsilon_{\mathrm{base}}$
in $> 90\%$ of T1 \\
H\ref{hyp:kruskal-wallis} & Scenario discrimination & Kruskal--Wallis $H$-test & $p < 0.01$ \\
H\ref{hyp:inter-layer-coupling} & Inter-layer coupling & Off-diagonal norm & $> 0$ (vs.\
permuted $\gamma$) \\
H\ref{hyp:severity-correlation} & Severity correlation & Spearman $\rho$ & $\rho > 0.5$,
$p < 0.01$ \\
H\ref{hyp:null-non-abelian} & Non-abelian structure & Permutation test & Exceeds $99$th
percentile of null \\
\bottomrule
\end{tabular}
\end{table}

\begin{remark}[Multiple Comparisons]
\label{rem:multiple-comparisons}
With six hypothesis tests, we apply Bonferroni correction to maintain a family-wise error
rate of $\alpha_{\mathrm{FWER}} = 0.01$. Each individual test is conducted at level
$\alpha / 6 \approx 0.0017$. For the Kruskal--Wallis test (Hypothesis~\ref{hyp:kruskal-wallis}),
post-hoc pairwise comparisons between scenario types are conducted using Dunn's test with
Holm--Bonferroni correction.
\end{remark}

\begin{remark}[Effect Sizes and Practical Significance]
\label{rem:effect-sizes}
Beyond statistical significance, we report effect sizes for all tests. For
Hypothesis~\ref{hyp:kruskal-wallis}, we report the epsilon-squared effect size
$\hat{\varepsilon}^2 = (H - k + 1) / (N - k)$, where $k = 3$ is the number of groups.
For Hypothesis~\ref{hyp:severity-correlation}, the Spearman $\rho$ itself serves as the
effect size. For block-specificity (Hypothesis~\ref{hyp:block-specificity}), we report
the mean dominant-layer ratio across conversations. These quantities enable readers to
assess the \emph{practical} significance of the curvature framework independent of sample
size.
\end{remark}

\begin{remark}[Reproducibility]
\label{rem:reproducibility}
All experimental stimuli, extraction code, and analysis scripts will be released upon
publication. The curvature computation pipeline (Definition~\ref{def:curvature-pipeline})
is implemented using TransformerLens~\cite{nanda2023transformerlens} for activation
extraction and NumPy/PyTorch for linear algebra operations. The head grouping function
$\gamma$ is estimated once on the full stimulus set and then fixed for all downstream
analyses, avoiding circularity.
\end{remark}
